\subsection{The finite soft coefficient in classical polylogarithms}
\label{sec:comp-ck-classical}
All weight-three terms in the preceding soft coefficient arise from
products of two ordinary logarithms.  This permits a classical
polylogarithmic endpoint expression for the complete coefficient,
without a GPL primitive operator.

For $z>0$ and real $\rho\ne0$, define
\begin{align}
 J(z,\rho)&=\ln^2z\ln|1-z/\rho|
 +2\ln z\operatorname{Re}\operatorname{Li}_2(z/\rho)
 -2\operatorname{Re}\operatorname{Li}_3(z/\rho),\\
 K(z,\rho)&=\ln z\ln|1-z/\rho|
 +\operatorname{Re}\operatorname{Li}_2(z/\rho).
\end{align}
The real parts denote the real boundary values on the standard
polylogarithm cut.  They enter because the original logarithms are real.
Their derivatives are listed below. At the degenerate values
$\rho=0,\infty$, the primitive is defined up to an additive constant,
which cancels in every normalized boundary difference. We use
\begin{align}
 \partial_z J(z,\rho)&=\frac{\ln^2z}{z-\rho},&
 J(z,0)&=\frac13\ln^3z,&J(z,\infty)&=0,\\
 \partial_z K(z,\rho)&=\frac{\ln z}{z-\rho},&
 K(z,0)&=\frac12\ln^2z.
\end{align}
The convention at infinity specifies a constant primitive with zero
derivative.  It will occur when a pole maps to infinity.

For a real letter $a$ outside the open physical interval, put
$A_a(z)=1-z/a$ and $\ell_a(z)=\ln A_a(z)$.  All factors used below
are positive for $0<z<Z$.  The normalized single-log primitive is
\begin{equation}
 \mathcal K_{a;c}(z)
 =K(A_a(z),A_a(c))-K(1,A_a(c)).
\end{equation}
For $a\ne b$ define
\begin{equation}
 \begin{aligned}
 \mathcal H_{a,b;c}(z)=\frac12\bigg[&
 J(A_a(z),A_a(c))+J(A_b(z),A_b(c))\\
 &-J\left(\frac{A_a(z)}{A_b(z)},
                 \frac{A_a(c)}{A_b(c)}\right)
 +J\left(\frac{A_a(z)}{A_b(z)},\frac ba\right)\bigg]
 -(z\to0).
 \end{aligned}
 \label{eq:comp-ck-classical-product}
\end{equation}
If $c=b$, the third term has its second argument at infinity and
vanishes.  If $a=b$, use
$\mathcal H_{a,a;c}(z)=J(A_a(z),A_a(c))-J(1,A_a(c))$.
These definitions obey
\begin{equation}
 \mathcal K'_{a;c}(z)=\frac{\ell_a(z)}{z-c},\qquad
 \mathcal H'_{a,b;c}(z)=\frac{\ell_a(z)\ell_b(z)}{z-c},\qquad
 \mathcal K_{a;c}(0)=\mathcal H_{a,b;c}(0)=0.
\end{equation}
To verify the product primitive, first use
$2\ln A\ln B=\ln^2A+\ln^2B-\ln^2(A/B)$.  Under
$v=A_a(z)/A_b(z)$, the rational differential transforms as
\begin{equation}
 \frac{\dd z}{z-c}
 =\frac{\dd v}{v-A_a(c)/A_b(c)}
  -\frac{\dd v}{v-b/a}.
\end{equation}
The second kernel is the image of infinity and is essential.
This identity derives both ratio terms in
\eqref{eq:comp-ck-classical-product}.  Positivity of $A_a,A_b$
justifies the logarithmic identity without additional multiples of
$\ii\pi$.  Subtraction at zero fixes every integration constant.

Here is a finite formula for including the rational coefficients.
For any rational $R(z)$ in the four soft kernels, write
\begin{equation}
 R=P_R'+\sum_c\frac{r_c(R)}{z-c},\qquad
 P_R=\sum_{n\geq0}\frac{p_nz^{n+1}}{n+1}
 -\sum_c\sum_{j\geq2}\frac{r_{c,j}}{(j-1)(z-c)^{j-1}}.
\end{equation}
The sums use its polynomial and partial-fraction coefficients.
They are finite algebraic derivatives of rational functions.
Let $\mathcal E[R](z)$ be the elementary primitive of $R$ normalized
to zero at the origin,
\begin{equation}
 \mathcal E[R](z)=P_R(z)-P_R(0)
 +\sum_cr_c(R)\ln\left|\frac{z-c}{-c}\right|.
\end{equation}
No pole here lies at the integration origin.  Define the following
explicit combinations of classical functions:
\begin{equation}
 \mathcal I_a[R]
 =P_R\ell_a+\sum_cr_c(R)\mathcal K_{a;c}
 -\mathcal E\!\left[\frac{P_R}{z-a}\right].
 \label{eq:comp-ck-classical-single}
\end{equation}
For the product, put $R_a=P_R/(z-a)$ and $R_b=P_R/(z-b)$.  Then
\begin{equation}
 \begin{aligned}
 \mathcal I_{a,b}[R]={}&P_R\ell_a\ell_b
 +\sum_cr_c(R)\mathcal H_{a,b;c}
 -P_{R_a}\ell_b-P_{R_b}\ell_a\\
 &-\sum_cr_c(R_a)\mathcal K_{b;c}
  -\sum_cr_c(R_b)\mathcal K_{a;c}\\
 &+\mathcal E\!\left[\frac{P_{R_a}}{z-b}
                         +\frac{P_{R_b}}{z-a}\right].
 \end{aligned}
 \label{eq:comp-ck-classical-double}
\end{equation}
These are nonrecursive formulas.  Direct differentiation gives
$\mathcal I_a'[R]=R\ell_a$ and
$\mathcal I_{a,b}'[R]=R\ell_a\ell_b$.
Their zero values at the origin follow from the normalized
$\mathcal E,\mathcal K,\mathcal H$ and $\ell_a(0)=0$.

Denote the weights of the three specified logarithm lists by
$b_a$ for $L_\beta$, $u_a$ for $L_1$ and $k_a$ for $L_k$.
The complete classical primitive is
\begin{equation}
 \begin{aligned}
 \Psi_{\mathrm{cl}}(z)={}&
 \sum_ab_a\mathcal I_a[\widetilde R_0]
 +\sum_{a,b}b_au_b\mathcal I_{a,b}[\widetilde R_1]\\
 &+\sum_{a,b}b_ak_b\mathcal I_{a,b}[\widetilde R_k]
 +\sum_ak_a\mathcal I_a[\widetilde R_f].
 \end{aligned}
\end{equation}
Each index ranges over its four explicitly given real letters.
Finally,
\begin{equation}
 \boxed{\mathcal C_K(s)=\lim_{z\to Z^-}\Psi_{\mathrm{cl}}(z)
 -\frac{\operatorname{artanh}\beta}{\beta}\mathcal C_0(s).}
 \label{eq:comp-ck-classical-complete}
\end{equation}
This expression uses only rational functions, real logarithms,
$\operatorname{Li}_2$ and $\operatorname{Li}_3$.  The endpoint limit
is taken after combining the terms.  It is the same finite limit
already established for the soft coefficient.

There is a further exact simplification before evaluation.  Each
logarithm list has weights summing to zero.  Hence the first two
$J$ terms in \eqref{eq:comp-ck-classical-product} cancel in both
weighted double sums, including their base-point constants.
Only the two ratio terms are needed there.  This cancellation
reduces the number of classical function evaluations without
changing the endpoint or its branch choice.
